\documentclass[11pt]{article}
\usepackage[T1]{fontenc}
\usepackage{lmodern,amsmath,amssymb,amsthm,microtype}
\usepackage[margin=1in]{geometry}
\usepackage[hidelinks]{hyperref}
\newtheorem{theorem}{Theorem}
\newtheorem{lemma}[theorem]{Lemma}
\newcommand{\F}{\mathbb F}
\newcommand{\eps}{\varepsilon}
\newcommand{\Th}{\Theta}
\title{A characteristic-two counterexample to d\'evissage\newline
for local deformation functors}
\author{Alec Kriebel\thanks{Independent researcher. ORCID:
\href{https://orcid.org/0009-0001-9320-500X}{0009-0001-9320-500X}.}}
\date{12 October 2026\quad (draft research note)}
\begin{document}
\maketitle
\begin{abstract}
We give a negative answer to the unrestricted fibre-product question for
local deformation functors posed by Byszewski. For the weakly ramified action
of the additive group of $\F_4$ given by $t\mapsto t/(1+at)$, and its subgroup
$\F_2$, restriction--induction cannot be formally smooth over any
prorepresentable recipient, regardless of its tangent dimension. The proof
combines the known involution deformation pathology with an exact calculation
showing that restriction has injective tangent map. The variant assuming
prorepresentability of the subgroup and quotient deformation functors is
untouched. The involution pathology and additive action are established prior
ingredients; the contribution is their use to obstruct the proposed
factorization.
\end{abstract}

\section{The unrestricted question and the counterexample}
Let $k$ be a perfect field of characteristic $p$ and $\mathcal C$ the category
of local Artinian $W(k)$-algebras with identified residue field $k$.
For a faithful local action $\rho:H\to\operatorname{Aut}_k k[[t]]$,
$D_H(A)$ is the set of lifts to $\operatorname{Aut}_A A[[t]]$, modulo
coordinate conjugation reducing to the identity. This is the unframed
category of \cite{OWR,Byszewski,BC}; nilpotent constant terms are permitted.
The residual action gives a unique point over $k$.
For $N\triangleleft G$ and $Q=G/N$, restriction and induction give a natural map
\[
 (\mathrm{res},\mathrm{ind}):D_G\longrightarrow D_N^Q\times D_Q,
\]
where $D_N^Q$ is the actual fixed subfunctor of $D_N$.
A prorepresentable functor has the form
\[
 h_R(A)=\operatorname{Hom}_{W(k),\mathrm{loc}}(R,A)
\]
for a complete Noetherian local $W(k)$-algebra with residue $k$.

The question in \cite[p.~3039]{OWR}, \cite[Question~6.4.4]{Thesis}, and
\cite[Question~4.4]{Byszewski} asks for a prorepresentable $F$, natural maps
$D_N^Q\to F\leftarrow D_Q$, and a formally smooth restriction--induction map
into the fibre product. The refined formulation prescribes
$TF=H^1(Q,\Th^\sharp/\Th^N)$, with
$\Th=\operatorname{Der}_k k[[t]]$ and
$\Th^\sharp=\operatorname{Der}_k(k[[t]]^N)$.
Neither question assumes $D_N$ or $D_Q$ prorepresentable. Those extra
hypotheses are stated separately in \cite[Theorem~3]{OWR} and
\cite[Remark~4.5, Theorem~4.9]{Byszewski}.

\begin{theorem}\label{thm:noF}
Put $k=\F_4=\F_2[u]/(u^2+u+1)$, $G=(k,+)$, $N=(\F_2,+)$, and
$\rho(a)(t)=t/(1+at)$. There are no prorepresentable $F$ and natural maps
$f:D_N^Q\to F$, $g:D_Q\to F$ for which restriction--induction factors as a
formally smooth map
\[
 D_G\longrightarrow D_N^Q\times_F D_Q.
\]
This remains impossible for every tangent dimension of $F$; in particular
it excludes the prescribed one-dimensional tangent space in this example.
\end{theorem}

The action is a standard weakly ramified additive action, not a new
construction. Direct substitution gives $\rho(a)\rho(b)=\rho(a+b)$, and the
$t^2$ coefficient distinguishes $a$. Every nonidentity element has conductor
one. The proof below uses actual deformation classes, without replacing
either component functor by a versal hull.

\section{A quotient tangent invisible to every recipient}
The $N$-quotient parameter and induced $Q$-action are
\[
 y=\frac{t^2}{1+t},\qquad k[[t]]^N=k[[y]],\qquad
 \sigma(y)=\rho(u)(y)=\frac{y}{1+y}.
\]
Here $t^2+yt+y=0$, $\rho(1)(t)=t+y$, and every series uniquely has the form
$A(y)+tB(y)$: the leading terms of $y^m$ and $ty^m$ have even and odd
$t$-orders, and successive leading-term subtraction converges.
Invariance forces $yB(y)=0$. The same basis argument for Laurent series
gives $k((t))^N=k((y))$. Finally $u(u+1)=1$ gives
\[
 \rho(u)(y)=\frac{t^2}{(1+ut)(1+(u+1)t)}
            =\frac{t^2}{1+t+t^2}=\frac{y}{1+y}.
\]
Thus $D_Q$ is the weakly ramified involution deformation functor.

\begin{lemma}\label{lem:recipient}
There is a nonzero $v\in D_Q(k[\eps]/\eps^2)$ whose image under every natural
map $D_Q\to h_R$ is the residue basepoint.
\end{lemma}
\begin{proof}
Let $B=k[\eps]/\eps^3$, $A=k[\eps]/\eps^2$, and let $j:B\to B$ fix $k$ and
send $\eps$ to $\eps+\eps^2$. These are objects and maps of $\mathcal C$ via
$W(k)\to k$. Define
\[
 \sigma_\eps(y)=\frac{y+\eps}{1+y},\qquad
 M_\eps=\begin{pmatrix}1&\eps\\1&1\end{pmatrix}.
\]
Its constant coefficient is nilpotent and its linear coefficient is a unit,
so substitution defines a continuous automorphism. Its matrix square is
$(1+\eps)I$, so it is an involution lifting $\sigma$.
Write $x=[\sigma_\eps]\in D_Q(B)$ and $v$ for its reduction to $A$.

The class $x$ is fixed by $j$. The identity-reducing coordinate change
$\chi(y)=(1+\eps)y+\eps$, represented by
$C=\bigl(\begin{smallmatrix}1+\eps&\eps\\0&1\end{smallmatrix}\bigr)$, satisfies
\[
 M_{\eps+\eps^2}C=(1+\eps)CM_\eps
\]
in $B$. Scalar matrices do not change the fractional-linear substitution,
so this gives the required conjugacy. Function composition reverses the
source's ring-substitution order; the equivalence is valid in either
conjugating direction. This is the known calculation of
\cite[Proposition~4.1]{BC}, reproduced here for completeness.

For a natural map $\eta:D_Q\to h_R$, write $\psi:R\to B$ for $\eta_B(x)$.
Naturality gives $j\circ\psi=\psi$. Since
\[
 B^j=k+k\eps^2,
\]
the composite $R\to B\to A$ lands in $k\subset A$. It is consequently the
residue map followed by $k\to A$, hence $\eta_A(v)$ is the basepoint. This
argument applies also to mixed-characteristic $R$ and arbitrary tangent
dimension, with no linearity assumption on $\eta$.

The class $v$ is nonzero. Every identity-reducing coordinate change over $A$
has the form $y+\eps h(y)$ with arbitrary $h\in k[[y]]$, including its
constant coefficient. Its first-order conjugacy perturbation is, up to sign,
$h(\sigma(y))-\sigma'(y)h(y)$, whose constant coefficient is zero.
But $\sigma_\eps-\sigma=\eps/(1+y)$ has constant perturbation coefficient
one. Thus no allowed conjugator trivializes $v$.
\end{proof}

The conclusion of Lemma~\ref{lem:recipient} is also an elementary consequence
of the previously published failure to preserve injections in
\cite[Proposition~6.1, p.~899]{BC}. There the nonzero involution tangent and
the basepoint become equal after the injection
$k[e]/e^2\hookrightarrow k[w]/w^3$, $e\mapsto w^2$.
Every $h_R$ preserves injections of target rings. Naturality therefore makes
their images equal already over $k[e]/e^2$. We claim no new recipient-collapse
phenomenon or new involution pathology; the direct cubic-ring proof fixes
the particular tangent class used here.

\section{The restriction tangent is injective}
Over $k$, every algebraic derivation of $k[[t]]$ is a coefficient derivation:
perfectness in characteristic two gives $h=a(t)^2+tb(t)^2$, so
$d(h)=d(t)b(t)^2=d(t)h'(t)$. Thus $\Th=k[[t]]\partial_t$, and these
derivations are continuous. There is no additional discontinuous tangent
module in this example.

Restriction to the $N$-fixed ring identifies
\[
 \Th^N\cong M:=yk[[y]]\partial_y.
\]
Indeed the unique extension of $b(y)\partial_y$ to the separable field
$k((t))$ has coefficient
\[
 a(t)=\frac{b(y)}{dy/dt}=b(y)\frac{(1+t)^2}{t^2},
 \qquad \frac{dy}{dt}=\frac{t^2}{(1+t)^2}.
\]
Uniqueness makes it $N$-invariant. It preserves $k[[t]]$ precisely when
$2\operatorname{ord}_y b-2\geq0$, namely $b\in yk[[y]]$.
Conversely an invariant derivation restricts and is recovered by this
formula. The field polynomial has derivative $y\ne0$, so separability
and uniqueness hold.

Conjugation by the quotient involution acts on coefficient functions by
\[
 T(f)(y)=(1+y)^2f\!\left(\frac{y}{1+y}\right),\qquad T^2=1.
\]
All group elements are involutions, so the inverse-conjugation convention
gives the same operator. Put $z=y^2/(1+y)$. The preceding fixed-ring basis
argument gives $k((y))^Q=k((z))$, and nonzero invariant Laurent series have
even $y$-valuation. For $f\in yk[[y]]$, $Tf=f$ is equivalent to invariance
of $f/y^2$. Its valuation is at least $-1$, hence by evenness at least zero.
It follows that
\[
 M^Q=y^2k[[z]]\partial_y,
 \qquad (T-1)(y h(z))=y^2h(z)\quad(h\in k[[z]]).
\]
The reverse containment of $(T-1)M$ in $M^Q$ follows from $(T-1)^2=0$.
The cyclic cocycle formula therefore gives
\[
 H^1(Q,\Th^N)=\frac{\ker(1+T)}{(T-1)M}=0.
\]
These are identities on complete formal series, not finite-jet calculations.

The first-order lifting calculation identifies $TD_H=H^1(H,\Th)$, and
restriction with cohomological restriction \cite[Proposition~2.4]{Byszewski}:
write a lift as $(\mathrm{id}+\eps c(g))\rho(g)$; group relations give a
cocycle, and coordinate conjugation changes it by a coboundary.
The vanishing above makes restriction injective. Explicitly, for a cocycle
whose restriction is zero in cohomology, subtract a coboundary so it vanishes
on $N$. The cocycle equation and commutativity of $G$ make its value at $u$
$N$-invariant; the relation $u+u=0$ gives $(1+T)c(u)=0$.
Write $c(u)=(T-1)m$ with $m\in\Th^N$ and subtract this coboundary too. It
preserves the zero $N$-values and kills $u$, hence kills the entire cocycle.
Thus
\[
 \ker(TD_G\longrightarrow TD_N)=0.
\]
Since $D_N^Q(A)\subset D_N(A)$ literally, restriction into the actual fixed
subfunctor also has zero tangent kernel.

For the prescribed space, $\Th^\sharp=k[[y]]\partial_y$ and
$\Th^\sharp/\Th^N\cong k$ has trivial $Q$-action (constant coefficients are
preserved by $T$). Hence $H^1(Q,\Th^\sharp/\Th^N)=k$, as asserted in the
theorem. Its dimension is not needed for the contradiction.

\section{Failure of formal smoothness and scope}
\begin{proof}[Proof of Theorem~\ref{thm:noF}]
Suppose such $F$, $f$ and $g$ exist; identify $F$ naturally with $h_R$.
At $A=k[\eps]/\eps^2$, naturality from $k$ sends the residual basepoint
$0_N\in D_N^Q(A)$ under $f$ to the residue basepoint of $F(A)$.
Lemma~\ref{lem:recipient} gives the same image for $g_A(v)$, with $v\ne0$.
Therefore $(0_N,v)$ belongs to $(D_N^Q\times_F D_Q)(A)$.
Its reduction is the unique residual point. The surjection $A\to k$ is
small: its kernel is annihilated by the maximal ideal of $A$.
Formal smoothness would lift this pair to $w\in D_G(A)$ with
\[
 \mathrm{res}(w)=0_N,\qquad\mathrm{ind}(w)=v.
\]
Restriction injectivity forces $w=0_G$. Naturality of induction then gives
$\mathrm{ind}(w)=0_Q$, contradicting $v\ne0$.
\end{proof}

The cubic ring is used only to constrain detection of a first-order tangent
by an arbitrary prorepresentable recipient. The obstruction to smoothness
already occurs on $A\to k$. No extension of $f$ from $D_N^Q$ to $D_N$ is
assumed, and no cohomological tangent formula for $f$ or $g$ is prescribed.
Nor is universality of a Klein-four deformation ring used.

This answers the unrestricted original question negatively and does not
contradict the conditional obstruction-space theorems cited above.
The involution component is in the exceptional nonprorepresentable case.
The variant imposing prorepresentability of $D_N$ and $D_Q$ remains
unaddressed; changing to coordinate automorphisms fixing the zero section
would likewise change the category and exclude the crucial nilpotent
translation. The introductory and refined formulations are one question,
not separate discoveries.

\paragraph{Priority, verification and review status.}
The known additive action, involution conjugacy, failure to preserve
injections, and tangent-cohomology framework are credited above. The
supplement records independent formal-series and functor audits, finite
exact checks and a bounded priority search. No matching prior negative
factorization was located in inspected primary bodies; this does not prove
exhaustive originality or that the question remained open until 2026.
Finite computations support explicit identities; the unrestricted formal
series and all-natural-map statements are proved in the text, without a
formal proof-assistant certificate.

AI tools were used extensively in solving, verifying, auditing sources, and
preparing this paper and supporting materials. This is an unrefereed research
note. Independent adversarial reviews are AI reviews; conventional human
peer review and journal acceptance are not claimed.

\begin{thebibliography}{9}
\bibitem{OWR} J.~Byszewski, \emph{D\'evissage for local deformation functors},
Oberwolfach Reports 54/2008, 3038--3040.
\href{https://ems.press/content/serial-article-files/46200}{Primary report}.
\bibitem{Thesis} J.~Byszewski, \emph{Cohomological aspects of equivariant deformation theory},
Ph.D. thesis, Utrecht University (2009),
Question~6.4.4.
\href{https://dspace.library.uu.nl/handle/1874/33749}{Repository record}.
\bibitem{Byszewski} J.~Byszewski,
\emph{Deformation functors of local actions},
\href{https://arxiv.org/abs/1112.0352v1}{arXiv:1112.0352v1} (2011).
\bibitem{BC} J.~Byszewski and G.~Cornelissen,
\emph{Which weakly ramified group actions admit a universal formal deformation?},
Ann. Inst. Fourier \textbf{59} (2009), 877--902.
\href{https://doi.org/10.5802/aif.2450}{doi:10.5802/aif.2450},
Propositions~4.1 and~6.1.
\end{thebibliography}
\end{document}
